\documentclass{article}

% if you need to pass options to natbib, use, e.g.:
%     \PassOptionsToPackage{numbers, compress}{natbib}
% before loading neurips_2026

% The authors should use one of these tracks.
% Before accepting by the NeurIPS conference, select one of the options below.
% 0. "default" for submission
\usepackage{neurips_2026}
% the "default" option is equal to the "main" option, which is used for the Main Track with double-blind reviewing.
% 1. "main" option is used for the Main Track
%  \usepackage[main]{neurips_2026}
% 2. "position" option is used for the Position Paper Track
%  \usepackage[position]{neurips_2026}
% 3. "eandd" option is used for the Evaluations & Datasets Track
 % \usepackage[eandd]{neurips_2026}
 % if you need to opt-in for a single-blind submission in the E&D track:
 %\usepackage[eandd, nonanonymous]{neurips_2026}
% 4. "creativeai" option is used for the Creative AI Track
%  \usepackage[creativeai]{neurips_2026}
% 5. "sglblindworkshop" option is used for the Workshop with single-blind reviewing
 % \usepackage[sglblindworkshop]{neurips_2026}
% 6. "dblblindworkshop" option is used for the Workshop with double-blind reviewing
%  \usepackage[dblblindworkshop]{neurips_2026}

% After being accepted, the authors should add "final" behind the track to compile a camera-ready version.
% 1. Main Track
 % \usepackage[main, final]{neurips_2026}
% 2. Position Paper Track
%  \usepackage[position, final]{neurips_2026}
% 3. Evaluations & Datasets Track
 % \usepackage[eandd, final]{neurips_2026}
% 4. Creative AI Track
%  \usepackage[creativeai, final]{neurips_2026}
% 5. Workshop with single-blind reviewing
%  \usepackage[sglblindworkshop, final]{neurips_2026}
% 6. Workshop with double-blind reviewing
%  \usepackage[dblblindworkshop, final]{neurips_2026}
% Note. For the workshop paper template, both \title{} and \workshoptitle{} are required, with the former indicating the paper title shown in the title and the latter indicating the workshop title displayed in the footnote.
% For workshops (5., 6.), the authors should add the name of the workshop, "\workshoptitle" command is used to set the workshop title.
% \workshoptitle{WORKSHOP TITLE}

% "preprint" option is used for arXiv or other preprint submissions
 % \usepackage[preprint]{neurips_2026}

% to avoid loading the natbib package, add option nonatbib:
%    \usepackage[nonatbib]{neurips_2026}

\usepackage[utf8]{inputenc} % allow utf-8 input
\usepackage[T1]{fontenc}    % use 8-bit T1 fonts
\usepackage{hyperref}       % hyperlinks
\usepackage{url}            % simple URL typesetting
\usepackage{booktabs}       % professional-quality tables
\usepackage{multirow}       % multi-row table cells
\usepackage{longtable}      % multi-page tables (appendix)
\usepackage{amsmath}        % math environments and \text
\usepackage{amsfonts}       % blackboard math symbols
\usepackage{nicefrac}       % compact symbols for 1/2, etc.
\usepackage{microtype}      % microtypography
\usepackage{xcolor}         % colors
\usepackage{graphicx}       % figures
\usepackage{subcaption}     % subfigures
\usepackage{algorithm}      % algorithm floats
\usepackage{algorithmic}    % algorithm pseudocode

\newcommand{\TanGCE}{TanGCE}
\newcommand{\AmbGCE}{AmbGCE}
\newcommand{\ye}[1]{\textcolor{purple}{(Yanai:) #1}}
\newcommand{\ma}[1]{\textcolor{blue}{(Matan:) #1}}
\newcommand{\yg}[1]{\textcolor{yellow}{(Yoav:) #1}}


% Note. For the workshop paper template, both \title{} and \workshoptitle{} are required, with the former indicating the paper title shown in the title and the latter indicating the workshop title displayed in the footnote. 
\title{Manifold-Aware Concept Erasure}


% The \author macro works with any number of authors. There are two commands
% used to separate the names and addresses of multiple authors: \And and \AND.
%
% Using \And between authors leaves it to LaTeX to determine where to break the
% lines. Using \AND forces a line break at that point. So, if LaTeX puts 3 of 4
% authors names on the first line, and the last on the second line, try using
% \AND instead of \And before the third author name.


\author{%
  David S.~Hippocampus\thanks{Use footnote for providing further information
    about author (webpage, alternative address)---\emph{not} for acknowledging
    funding agencies.} \\
  Department of Computer Science\\
  Cranberry-Lemon University\\
  Pittsburgh, PA 15213 \\
  \texttt{hippo@cs.cranberry-lemon.edu} \\
  % examples of more authors
  % \And
  % Coauthor \\
  % Affiliation \\
  % Address \\
  % \texttt{email} \\
  % \AND
  % Coauthor \\
  % Affiliation \\
  % Address \\
  % \texttt{email} \\
  % \And
  % Coauthor \\
  % Affiliation \\
  % Address \\
  % \texttt{email} \\
  % \And
  % Coauthor \\
  % Affiliation \\
  % Address \\
  % \texttt{email} \\
}


\begin{document}

% Manuscript entrypoint: include modular section files from sections/
\maketitle

% Abstract input (use without .tex suffix)
% Abstract
\section*{Abstract}

Concept erasure should remove a target concept without damaging the rest of the representation, but nonlinear concepts remain recoverable after standard linear interventions and unconstrained ambient-space edits can distort representations. We propose \TanGCE{}, a new manifold-aware erasure method that estimates local tangent spaces with neighborhood PCA and restricts erasure updates to that local support. The core idea is simple: if natural representations lie on a low-dimensional manifold inside the ambient space, then effective interventions should follow that manifold rather than push representations away from it. Across our completed text and vision benchmarks, \TanGCE{} delivers the strongest overall nonlinear erasure results in the study: it consistently beats the prior learned baselines, including \AmbGCE{}, and avoids the severe control damage incurred by distortion-based baselines. More broadly, these results suggest that manifold-aware editing is a useful principle not only for concept removal, but for representation editing and optimization in high-dimensional models.


% Main sections: wired to separate files so content can be edited independently
% Introduction
\section{Introduction}

Effective concept erasure must do two things at once: remove the target concept and preserve the rest of the representation. This is especially difficult for nonlinear concepts: an edit may defeat one trained probe while leaving behind signal that a freshly trained nonlinear classifier can still recover.
Recent studies have primarily examined linear concept erasure \citep{ravfogel-etal-2020-null, belrose-etal-2023-leace} \ma{cite}, providing strong theoretical guarantees under various assumptions about when a concept has been removed from a representation \ye{why do the theoretical guarantees matter here? remove}. Yet modern neural networks are fundamentally nonlinear, and the concepts encoded in their activations are therefore often nonlinear as well. This leaves nonlinear concept erasure as an important open problem for practical applications, including bias mitigation, fairness, and decision-making systems involving non-subjective concepts.
\ye{i would start with this last sentence: we want to make models better/more fair for these different applications. one way to do this is to erase the concepts from representations. however, concepts are often non-linear, and while recent work has focused on linear concepts, this doesn't work well when concepts are non-linear}\ma{The issue is that we target static activations in our experiment, and dont have experiments on model predictions using those modified activations}


% as emphasized by Elazar and Goldberg's cautionary analysis of post-hoc recovery after adversarial removal \citep{elazar-goldberg-2018-adversarial}. For erasure to be useful, it must therefore achieve two goals simultaneously: high \emph{erasure precision} on the target concept and high \emph{surgicality} on control concepts. \ye{the focus here on our paper will beg the question of why we didn't compare to that paper's method}
% \ye{but also, it's unclear why is this the first thing to mention in the introduction. this is an important detail for the evaluation, but it's not the core idea of this paper}\ye{the first sentence and a half are good. continue from there}\ye{last sentence is repetitive of the first one. remove.}

We approach this problem through a local manifold view of representation space. Hidden states live in a high-dimensional ambient space \ye{this term isn't clear/known. you should explain what you mean by that, and use italics the first time you use it}, but the representations used by natural inputs need not fill that space. Instead, they may concentrate near a much smaller local support. Under this view, the main geometric risk is not simply that an edit is global, but that an ambient-space update can move representations off the natural support where the concept actually lives.

This suggests a different design principle for nonlinear erasure. Rather than editing along the full ambient gradient \ye{again, unclear}, we estimate a local tangent approximation around each sample and retain only the component of the update supported by that local geometry. In practice, we approximate the tangent space with local PCA/SVD \ye{which one is it? are these two options? can we just say matrix decomposition?} on neighborhood secants. This is a first-order approximation rather than a theorem about hidden-state geometry, but it yields a concrete and testable erasure method.

Our contributions are twofold:
\begin{enumerate}
    \item \textbf{A manifold-motivated formulation of nonlinear concept erasure.} We frame nonlinear erasure through a local manifold hypothesis: the relevant question is not only whether representations are low-rank, but whether editing should respect the local support rather than act directly in ambient space. \ye{why is this not part of the second point?}

    \item \textbf{A tangent-constrained nonlinear erasure method.} In Section~\ref{sec:method}, we introduce \TanGCE{}, which operationalizes this view by projecting ambient erasure updates onto locally estimated tangent spaces so that interventions follow the first-order geometry supported by nearby representations. \ye{it's unclear what you method does. you're using all these terms that you don't explain and just sound fancy. the beauty is in simplicity}
    \ye{also, you didn't mention at all the part about using knn of other inputs to help you locate the manifold. it's very implicit in the writing, and you should be explicit about it}
\end{enumerate}

We therefore evaluate with a recoverability-oriented protocol: after erasure, we retrain a fresh nonlinear probe and iterate the intervention to remove residual signal. Across our completed text and vision evaluations \ye{first time you mention those. it sounds like you already mentioned them. you should have a paragraph stating you run extensive experiments and evaluation on a variety of datasets and modalities, as well as models (of diff families and sizes)}, \TanGCE{} is the stronger nonlinear remover than its ambient ablation \ye{unclear what ambient ablation is} and the earlier methods we compare against. The geometry results \ye{what's that?} support the manifold motivation in a bounded form: representations show local low-rank structure, but ambient edits can still leave the natural support. Section~\ref{sec:related} positions the method against prior erasure work, Section~\ref{sec:method} gives the method and modeling view, Section~\ref{sec:evaluation} gives the experimental setup, and Section~\ref{sec:results} evaluates both erasure performance and the geometry motivating local tangent constraints.

% Related work
\section{Related Work}
\label{sec:related}

Concept erasure methods differ mainly in the function class against which they erase and in the geometry they assume. Linear methods such as INLP \citep{ravfogel-etal-2020-null} and LEACE \citep{belrose-etal-2023-leace} target linearly decodable concept information. INLP removes successively exposed linear directions by iterative nullspace projection, while LEACE gives the strongest closed-form linear baseline in our comparisons. These methods are essential reference points, but they do not directly address the case where a fresh nonlinear probe can still recover the concept after linear erasure.

That recoverability concern is central to our framing. Elazar and Goldberg \citep{elazar-goldberg-2018-adversarial} showed that apparent success against one learned remover or classifier does not imply that the protected information is gone for a newly trained downstream model. Amnesic probing \citep{elazar-etal-2021-amnesic} and related intervention-based analyses adopt the same conservative lesson: the meaningful question is what remains recoverable after the intervention, not only what the original scorer no longer captures.

Two nonlinear erasure methods anchor our comparisons. IGBP \citep{iskander-etal-2023-shielded} is the closest prior baseline in mechanism: it iteratively trains a nonlinear probe and retracts representations toward the current decision boundary. Obliviator \citep{akbari-etal-2025-obliviator} is the strongest published baseline in headline numbers: it formulates erasure in RKHS with a Hilbert--Schmidt independence criterion (HSIC) penalty between the representation and the target attribute, and gradually morphs the feature space via random Fourier feature eigenvalue projection (RFF-EVP). Both methods are supervised on the target concept and unsupervised on control concepts (the same regime as our method); they differ from \TanGCE{} in how the update is constructed, not in what supervision they use. We report INLP, IGBP, and Obliviator under the same trajectory-probe protocol as every other method in Tab.~\ref{tab:main_results_unified}. Our paper differs in two ways. First, \AmbGCE{} is introduced only as the direct ambient-space ablation that keeps the iterative nonlinear scorer loop but removes the manifold constraint. Second, \TanGCE{} is the paper's main method: it projects the erasure update onto a locally estimated tangent space, making the geometric hypothesis itself the main point of comparison.

% Problem setup and geometric motivation
\section{Problem setup and geometric motivation}
\label{sec:notation}

Let $\mathrm{enc}_\ell : \Sigma^* \to \mathbb{R}^d$ be the hidden state at layer $\ell$ of a language model over vocabulary $\Sigma$. For each sample $i$ we have a representation $\mathbf{x}_i = \mathrm{enc}_\ell(s_i) \in \mathbb{R}^d$, a binary concept label $y_i \in \{0,1\}$, and an optional control label $z_i$. The dataset is $\mathcal{D} = \{(\mathbf{x}_i, y_i, z_i)\}_{i=1}^N$ with class subsets $\mathcal{D}_c = \{i : y_i = c\}$. A concept scorer $f_\theta : \mathbb{R}^d \to \mathbb{R}$ is a 2-layer MLP with input gradient $\nabla f_\theta(\mathbf{x}_i)$.

\paragraph{Local manifold hypothesis.} We adopt the modeling assumption that the representations $\{\mathbf{x}_i\}$ lie near a locally low-dimensional support $\mathcal{M} \subset \mathbb{R}^d$. Its operational consequence is that we can approximate the neighborhood of any $\mathbf{x} \in \mathcal{M}$ by its tangent space $T_{\mathbf{x}}\mathcal{M}$ \citep{tu-2011-manifolds} and apply linear-algebraic operations (SVD, projection, anisotropic scaling) inside it. The tangent space is not observed directly; we estimate it from nearby representations by local SVD on the secant matrix (§\ref{sec:method}). For an ambient vector $\mathbf{g} \in \mathbb{R}^d$, the decomposition $\mathbf{g} = \mathbf{g}^{\top} + \mathbf{g}^{\perp}$ with $\mathbf{g}^{\top} \in T_{\mathbf{x}}\mathcal{M}$ separates the part of $\mathbf{g}$ supported by the local sample structure from the part orthogonal to it. \TanGCE{} acts on the tangent component only; \AmbGCE{} retains both. Formal coordinate-chart definitions and the differential $d(\varphi^{-1})_0$ identification with $T_{\mathbf{x}}\mathcal{M}$ are deferred to App.~\ref{sec:appendix_geometry}.

\paragraph{Recoverability.} Concept removal is judged only after training a \emph{fresh} probe on the edited representations \citep{elazar-goldberg-2018-adversarial}: erasure succeeds if a newly trained nonlinear classifier fails to recover the concept from $\tilde{\mathbf{x}}$, not merely if the old scorer is fooled.

% Method — Implementation Details
% (Formal definitions are in Section~\ref{sec:notation}; this section describes the iterative protocol.)

\section{Implementation}
\label{sec:method}

We implement GCE using an iterative protocol with probe recycling. Each round of erasure begins by training a fresh concept scorer on the current (possibly already partially erased) representations, computing per-sample gradients, and applying the erasure step (Eq.~\ref{eq:ambient_erase} for Ambient GCE, Eq.~\ref{eq:tangent_erase} for TanGCE). The tangent estimator is refit in every round to track the evolving manifold geometry as representations shift.

Two diagnostics verify tangent quality throughout the iterative process. \emph{Tangent Gradient Capture} (TGC) measures the fraction of the leakage gradient that lies within the estimated tangent space. \emph{Tangent Neighbor Residual} (TNR) measures the fraction of local secant vectors not explained by the tangent basis. High TGC and low TNR confirm that the tangent estimation remains geometrically sound.

% Experimental setup
\section{Experimental Setup}
\label{sec:evaluation}

We evaluate each erasure method by training all probes from scratch on the intervened representations, so that success cannot be attributed to leakage from the erasure scorer itself.

\subsection{Evaluation Protocol}

\paragraph{Fresh-probe evaluation.}
The precision metric is intentionally strict: after every erasure method finishes, we train a \emph{new} probe on the intervened representations. If a freshly trained probe can still recover the concept, then the concept was not actually removed in the sense relevant for downstream use.

\paragraph{Shared validation protocol.}
Most gradient-based methods share a common erasure strength parameter $\lambda$. For \AmbGCE{} and \TanGCE{}, $\lambda$ scales the ambient or tangent-projected update. Because \TanGCE{}'s tangent projection attenuates the effective step size by $\|\mathbf{u}_i^{\mathrm{tan}}\|^2 \leq 1$, a fixed $\lambda$ across methods would systematically under-erase for the tangent-constrained update. Rather than manually choosing different values, we adopt validation-based hyperparameter selection: all $\lambda$-based methods search the same grid and select the $\lambda^*$ that minimizes validation nonlinear probe accuracy. The main comparison (Table~\ref{tab:full_comparison}) and vision benchmarks use $\lambda \in \{0.5, 1.0, 2.0, 4.0, 8.0\}$; the multi-model generalization experiments use the extended grid $\lambda \in \{0.25, 0.5, 1.0, 2.0, 4.0, 8.0, 16.0, 32.0\}$, which we found necessary for larger models and deeper layers where optimal $\lambda^*$ can exceed $8.0$. IGBP has no explicit $\lambda$ parameter---its step size is determined by each sample's distance to the decision boundary---so it runs until convergence.

\paragraph{Training protocol.}
All MLP probes---the shared \emph{concept scorer} used as the erasure direction source, the \emph{validation probe} for $\lambda$ HPO, and the \emph{evaluation probe} retrained on intervened representations---share the same architecture: two-layer MLP with hidden dimension 256, AdamW (lr~$=10^{-3}$, weight decay $10^{-4}$), batch size 512. For language model experiments, all three probes use 1500 gradient steps. For CelebA, the concept scorer uses 1000 steps and the HPO and evaluation probes use 1200 steps; the shorter budget suffices for the lower-dimensional CLIP embeddings. Using identical probe budgets for HPO and evaluation within each modality ensures that the lambda selected on validation is directly comparable to the final test metric, with no ranking distortion from mismatched model capacity. The tangent estimator uses $k{=}50$ nearest neighbors for local PCA; the retained tangent rank per sample is determined by a retained-variance threshold applied to the local eigenspectrum. IGBP runs for up to 35 rounds; each round trains its probe for 10 epochs with AdamW lr~$=2{\times}10^{-4}$, batch size 256, and stops early when validation accuracy falls to majority-class level. Linear baselines (LEACE, INLP, Linear Projection) use $\ell_2$-regularized logistic regression with balanced class weights. The control-task predictor is also an $\ell_2$ logistic regression trained on clean representations and evaluated on the intervened test set without retraining.

\paragraph{Domain-specific controls.}
For sycophancy we use answer matching as the control concept; for Bias in Bios we use profession; for safety we use helpfulness; and for CelebA we use selected low-correlation facial attributes (5 per concept). This lets the evaluation distinguish targeted concept removal from broad representational damage.

\paragraph{Vision benchmarks.}
For cross-modal evaluation we use CelebA~\citep{liu-etal-2015-celeba}, which provides 202{,}599 celebrity images with 40 binary facial attributes; we use the official train/val/test split. We extract pooled CLIP ViT-B/32 image embeddings (768-dim) and apply all erasure methods to the resulting representation matrix.

\subsection{Metrics}

\paragraph{Erasure precision.}
Let $h^{\mathrm{NL}}$ be a nonlinear probe and $h^{\mathrm{L}}$ a linear probe trained on $\tilde{\mathbf{x}}_i$ to predict the target concept $y_i$. We report their test accuracies on the intervened representations:
\begin{equation}
\label{eq:precision_metric}
\mathrm{Prec}_{\mathrm{erase}}^{\mathrm{NL}} = \mathrm{Acc}(h^{\mathrm{NL}}(\tilde{\mathbf{X}}), y),
\qquad
\mathrm{Prec}_{\mathrm{erase}}^{\mathrm{L}} = \mathrm{Acc}(h^{\mathrm{L}}(\tilde{\mathbf{X}}), y).
\end{equation}
Lower values indicate stronger erasure. Because the target concepts in this paper are not assumed to be linearly encoded, the retrained nonlinear probe is the primary precision metric and the linear probe is a secondary diagnostic. For interpretation, the majority-class predictor provides a trivial lower-bound reference for what it means to erase by collapsing the target signal entirely rather than preserving useful structure.

\paragraph{Surgicality on control concepts.}
Let $q$ be a predictor for an annotated control label $c$ (e.g., profession for gender erasure, helpfulness for safety, or selected low-correlation attributes in CelebA). We report the change in control performance relative to the clean representations:
\begin{equation}
\label{eq:surgicality_metric}
\Delta_{\mathrm{ctrl}} = \mathrm{Acc}(q(\tilde{\mathbf{X}}), c) - \mathrm{Acc}(q(\mathbf{X}), c).
\end{equation}
Smaller $|\Delta_{\mathrm{ctrl}}|$ means the erasure is more surgical. In CelebA, where each target attribute is paired with multiple controls, we report the mean $\Delta_{\mathrm{ctrl}}$ across the selected controls.

Together, these metrics ask whether an erasure method removes the target concept while preserving unrelated control concepts.

% Experiments
\section{Experiments}
\label{sec:experiments}

Our experimental evaluation focuses on characterizing the performance of \texttt{gce\_tangent} against \texttt{gce\_ambient} across three distinct concept domains. We prioritize nonlinear concepts where manifold geometry is expected to play a critical role, while maintaining a diagnostic baseline for linear comparison.

\subsection{Experimental Domains and Ordering}
We evaluate our methods across three primary blocks, ordered by their relevance to nonlinear manifold erasure:
\begin{enumerate}
    \item \textbf{Sycophancy (Primary Nonlinear Case):} Our central focus is on removing agreement bias in model responses. This represents a critical challenge for model alignment where concept encoding is expected to be locally curved.
    \item \textbf{Safety (Secondary Case Study):} We evaluate the surgicality of erasure on content safety using the PKU-SafeRLHF dataset, testing the preservation of helpfulness alongside safety leakage.
    \item \textbf{Bias in Bios (Diagnostic Baseline):} We use gender and profession classification as a replication anchor. This setting provides a linear comparison point to contrast against the nonlinear gains of tangent-constrained erasure.
\end{enumerate}

\subsection{Model and Implementation}
All experiments are conducted on Qwen2.5-0.5B (24 decoder layers, hidden dimension 896). The current paper-locked comparison set implements five primary erasure methods: linear projection, INLP, LEACE, \texttt{gce\_ambient}, and \texttt{gce\_tangent}. LEACE is included as a matched linear baseline in the layer-8 inventory reported below. For \texttt{gce\_tangent}, we estimate local tangent spaces with PCA-based neighborhoods so that the erasure direction is constrained to a local on-manifold approximation rather than to the full ambient space.

\begin{table}[t]
\centering
\caption{Comparison of concept erasure methods by family and role in this work.}
\label{tab:method_comparison_scope}
\scriptsize
\begin{tabular}{llp{0.48\linewidth}}
\toprule
Method & Family & Role in Paper \\
\midrule
Linear Projection & Linear & Global baseline \\
INLP & Linear & Iterative linear baseline \\
LEACE & Linear & Strong matched linear baseline \\
IGBP & Nonlinear & Closest prior nonlinear baseline \\
\AmbGCE{} & Nonlinear & Ambient ablation of \TanGCE{} \\
\TanGCE{} & Nonlinear & Main contribution \\
\bottomrule
\end{tabular}
\end{table}


\subsection{Evaluation Protocol}
All leakage probes are trained from scratch on erased representations to ensure no information leaks from the erasure scorers themselves. We employ three complementary evaluation axes:

\paragraph{Probe-based leakage.} A retrained nonlinear MLP probe measures whether concept information remains accessible to a standard function class. We report both linear and nonlinear test accuracy; the gap between them indicates the degree to which concept encoding is nonlinear.

\paragraph{Mutual information via $k$-NN estimation.}
Probe accuracy measures whether a \emph{specific function class} can recover the concept; a nonlinear MLP probe might miss information that a more complex architecture could extract. Mutual information $I(\mathbf{X}; y)$ captures \emph{all} shared information between representations and labels, providing a probe-free, architecture-agnostic bound on recoverability.

\emph{Mental model.} Mutual information decomposes as $I(\mathbf{X}; y) = H(y) - H(y \mid \mathbf{X})$: the entropy of the concept label minus the residual uncertainty after observing the representation. When erasure is effective, $H(y \mid \mathbf{X})$ rises toward $H(y)$ (the label becomes unpredictable from the representation), and MI drops toward zero. For continuous $\mathbf{X}$ and discrete $y$, we estimate MI using the Ross/Gao $k$-nearest-neighbor estimator:
\begin{equation}
\label{eq:knn_mi}
I(\mathbf{X}; y) = \psi(k) - \bigl\langle \psi(m_i) \bigr\rangle + \psi(N) - \bigl\langle \psi(N_{c(i)}) \bigr\rangle,
\end{equation}
where $\psi$ is the digamma function and the procedure is:
\begin{enumerate}
    \item For each point $\mathbf{x}_i$, find its $k$-th nearest neighbor \emph{within its class} $c(i)$; call that distance $\varepsilon_i$.
    \item Count $m_i$: the number of \emph{all} points (any class) within distance $\varepsilon_i$ of $\mathbf{x}_i$.
    \item $N_{c(i)}$ is the total number of points sharing the same label as point $i$; $N$ is the total sample size.
\end{enumerate}

\emph{Intuition.} When concept information is encoded in $\mathbf{X}$, same-class points cluster together: $\varepsilon_i$ is small, $m_i \approx k$ (mostly same-class neighbors within that radius), and MI is high. After effective erasure, class labels become geometrically mixed: to reach $k$ same-class neighbors one must search a large $\varepsilon_i$, sweeping in many cross-class points so that $m_i \gg k$, and MI drops.

\emph{Adaptive resolution and local geometry.} A key property of the $k$-NN estimator is that it sets its measurement scale \emph{adaptively}: for each point, the radius $\varepsilon_i$ is determined by the local data density, not by a global bandwidth or histogram bin width. In dense regions of the manifold, $\varepsilon_i$ is small; in sparse regions, it is large. The digamma terms $\psi(m_i)$ implicitly encode the local density ratio between the class-conditional and marginal distributions without ever estimating densities explicitly in $\mathbb{R}^d$---the estimator works entirely with neighbor counts, which are topological quantities that naturally handle varying density and curvature.

This makes the estimator well-suited for data on a low-dimensional manifold embedded in $\mathbb{R}^d$. A parametric density estimator (e.g., Gaussian kernel) operating in the ambient space would be misled by the high-dimensional volume: two points close in $\mathbb{R}^d$ may be far apart along the manifold, and vice versa. The $k$-NN estimator sidesteps this because it only examines local neighborhoods---points that are actually nearby along the data distribution. While it still uses ambient Euclidean distance (not geodesic distance), the locality of $k$-NN queries means that for small $k$ relative to the dataset, the Euclidean and geodesic neighborhoods approximately coincide.

\emph{Connection to TanGCE.} Both the MI estimator and TanGCE operate through local neighborhood structure rather than global ambient-space assumptions. TanGCE explicitly projects onto the tangent space estimated from $k$-nearest neighbors; the MI estimator implicitly adapts to the manifold by measuring information through local ball neighborhoods. Neither assumes the data fills $\mathbb{R}^d$ uniformly. This shared reliance on local structure creates a coherent geometric framework: the same manifold properties that make tangent-projected erasure more precise also make $k$-NN MI estimation more faithful than ambient-space alternatives.

We compute MI for both concept labels (should decrease after erasure) and control labels (should remain unchanged), yielding a complete picture of erasure selectivity.

\paragraph{Unsupervised SAE diagnostics.} We train TopK sparse autoencoders ($k{=}64$) on a general-purpose WikiText corpus---independent of any concept labels---then measure concept feature removal rate, collateral damage rate, and feature redistribution score (Section~\ref{sec:results}). These metrics require no control-task annotations and detect failure modes invisible to probe-based evaluation, such as concept information being redistributed across features rather than removed.

\subsection{Geometric Foundation (Appendix Support)}
While the main text focuses on the performance of \texttt{gce\_tangent}, we also provide geometric diagnostics to ground the mechanism story. These include broad geometry sweeps across multiple datasets and layers, subspace projection sensitivity (Q1), perturbation-based subspace sensitivity tests (Q2), and representation connectivity analyses (Q3). These appendix-scale materials support the low-dimensional and nonlinear-geometry motivation for tangent projection, but they should not be read as a standalone causal proof that off-manifold displacement fully explains the observed utility tradeoffs.

% Results
\section{Results}
\label{sec:results}

We evaluate seven erasure methods across three concept domains---sycophancy, gender (Bias in Bios), and safety---at layer~8 of Qwen2.5-0.5B. All methods except TanGCE use $\lambda{=}0.5$; TanGCE uses $\lambda{=}8.0$ to compensate for the naturally attenuated projection norms (see Section~\ref{sec:lambda_sweep}). We measure concept leakage with a retrained nonlinear MLP probe, mutual information via the Ross/Gao $k$-NN estimator in the original representation space, and unsupervised SAE-based metrics: concept feature removal rate, collateral damage rate, and feature redistribution score.

\begin{table}[t]
\centering
\caption{Full erasure method comparison at layer 8 across three concept domains. All methods use $\lambda{=}0.5$ except TanGCE which uses $\lambda{=}8.0$ to compensate for attenuated projection norms (Section~\ref{sec:lambda_sweep}). L/NL Leak: retrained linear logistic-regression / nonlinear MLP probe accuracy (lower = better concept removal). $\Delta$Ctrl: downstream control task accuracy change relative to clean representations (closer to zero = more surgical). Dmg\%: fraction of non-concept SAE features with $>{10}\%$ activation change (lower = less collateral damage). Rdst\%: fraction of non-concept SAE features that become concept-correlated post-erasure (lower = less information reshuffling). Model: Qwen2.5-0.5B. SAE: $64\times$ TopK, $k{=}64$, disruption threshold $\delta{>}10\%$. Control tasks: answer matching (sycophancy), profession classification (gender), helpfulness (safety).}
\label{tab:full_comparison}
\small
\setlength{\tabcolsep}{3pt}
\begin{tabular}{l|ccccc|ccccc|ccccc}
\toprule
 & \multicolumn{5}{c|}{Sycophancy} & \multicolumn{5}{c|}{Gender (Bias in Bios)} & \multicolumn{5}{c}{Safety} \\
Method & L$\downarrow$ & NL$\downarrow$ & $\Delta$Ctrl & Dmg & Rdst & L$\downarrow$ & NL$\downarrow$ & $\Delta$Ctrl & Dmg & Rdst & L$\downarrow$ & NL$\downarrow$ & $\Delta$Ctrl & Dmg & Rdst \\
\midrule
Random Global       & .685 & .748 & .000 & 0.0  & 0.0  & .719 & .758 & .000  & 0.8  & 0.0  & .798 & .805 & +.003 & 0.0  & 0.0  \\
Random Per-Samp.    & .621 & .745 & .000 & 0.0  & 0.0  & .723 & .743 & +.003 & 0.0  & 0.0  & .784 & .799 & $-.010$ & 0.0  & 0.0  \\
Linear Proj.        & .683 & .740 & .000 & 5.3  & 0.0  & .722 & .747 & +.001 & 0.3  & 0.0  & .807 & .800 & $-.003$ & 0.0  & 0.0  \\
INLP                & .683 & .734 & .000 & 5.3  & 0.0  & .722 & .737 & +.001 & 0.3  & 0.0  & .807 & .795 & $-.004$ & 0.0  & 0.0  \\
LEACE               & .681 & .681 & .000 & 0.0  & 0.0  & .719 & .739 & $-.002$ & 0.8  & 0.0  & .801 & .785 & $-.015$ & 0.0  & 0.0  \\
Ambient GCE         & .714 & .703 & .000 & 0.0  & 0.0  & .663 & .701 & .000  & 0.2  & 0.0  & .780 & .778 & $-.040$ & 0.0  & 3.3  \\
\textbf{TanGCE} & \textbf{.628} & \textbf{.602} & $-.015$ & 18.4 & \textbf{0.0} & \textbf{.603} & \textbf{.618} & $-.008$ & 9.8 & \textbf{0.0} & \textbf{.654} & \textbf{.663} & +.039 & 21.3 & 1.6 \\
\bottomrule
\end{tabular}
\end{table}

\begin{table}[t]
\centering
\caption{Mutual information and fidelity metrics (layer 8). MI$_c$: concept mutual information in original representation space via $k$-NN estimator (lower = better erasure). MI$_\text{ctrl}$: control-label mutual information (higher = better preservation of task-relevant information). Clean baselines: sycophancy MI$_c{=}0.189$, MI$_\text{ctrl}{=}1.157$; gender MI$_c{=}0.074$, MI$_\text{ctrl}{=}0.903$; safety MI$_c{=}0.160$, MI$_\text{ctrl}{=}0.077$. $L_2$: mean $\ell_2$ displacement from clean representations.}
\label{tab:full_comparison_detail}
\small
\setlength{\tabcolsep}{3pt}
\begin{tabular}{l|cccc|cccc|cccc}
\toprule
 & \multicolumn{4}{c|}{Sycophancy} & \multicolumn{4}{c|}{Gender (Bias in Bios)} & \multicolumn{4}{c}{Safety} \\
Method & MI$_c\!\downarrow$ & MI$_\text{ctrl}$ & $L_2$ & Ctrl & MI$_c\!\downarrow$ & MI$_\text{ctrl}$ & $L_2$ & Ctrl & MI$_c\!\downarrow$ & MI$_\text{ctrl}$ & $L_2$ & Ctrl \\
\midrule
Random Gl.    & .191 & 1.155 & 0.02 & 1.000 & .075 & .902 & 0.10 & .538 & .164 & .075 & 0.06 & .565 \\
Linear Pr.    & .191 & 1.155 & 0.11 & 1.000 & .074 & .902 & 0.06 & .540 & .161 & .074 & 0.01 & .560 \\
INLP          & .195 & 1.158 & 0.11 & 1.000 & .076 & .902 & 0.06 & .540 & .161 & .075 & 0.01 & .559 \\
LEACE         & .198 & 1.159 & 0.07 & 1.000 & .075 & .901 & 0.20 & .537 & .127 & .070 & 0.15 & .548 \\
Amb.~GCE      & .184 & 1.158 & 0.02 & 1.000 & .077 & .902 & 0.08 & .538 & .157 & .075 & 0.04 & .523 \\
\textbf{TanGCE} & \textbf{.113} & 1.158 & 0.86 & .985 & \textbf{.073} & \textbf{.904} & 0.41 & .531 & \textbf{.145} & .069 & 0.65 & .602 \\
\bottomrule
\end{tabular}
\end{table}


\subsection{TanGCE Achieves Lowest Leakage Across All Domains}

Table~\ref{tab:full_comparison} shows that TanGCE achieves the lowest nonlinear leakage across all three domains: 0.602 on sycophancy (vs.\ 0.681 for LEACE, 0.703 for Ambient GCE), 0.618 on gender (vs.\ 0.701 for Ambient GCE), and 0.663 on safety (vs.\ 0.778 for Ambient GCE). Linear methods (Linear Projection, INLP) barely reduce leakage below random baselines, confirming that concept information is predominantly encoded nonlinearly.

\subsection{Control Task Preservation}

Table~\ref{tab:full_comparison} reports $\Delta$Ctrl, the change in downstream control task accuracy relative to clean representations. For sycophancy (answer matching), TanGCE incurs a small $-1.5$pp drop; all other methods preserve perfect accuracy. For gender (profession classification), TanGCE drops only $-0.8$pp (0.531 vs.\ 0.539 clean). On safety, TanGCE \emph{improves} helpfulness ($+3.9$pp), suggesting that removing safety-concept directions frees capacity for the helpfulness task.

Table~\ref{tab:full_comparison_detail} complements this with mutual information: TanGCE preserves control MI across all domains (sycophancy: 1.158 vs.\ 1.157 clean; gender: 0.904 vs.\ 0.903 clean). On gender, TanGCE is the only method that slightly \emph{increases} control MI, indicating zero collateral damage to task-relevant information.

\subsection{Near-Zero Redistribution}

TanGCE maintains near-zero redistribution across all concepts: 1.4\% on sycophancy, 0.0\% on gender, and 4.1\% on safety. In contrast, Ambient GCE shows nonzero redistribution even at $\lambda{=}0.5$, and this grows dramatically at higher $\lambda$ (Section~\ref{sec:lambda_sweep}).

\paragraph{Damage threshold sensitivity.}
\begin{table}[t]
\centering
\caption{Effect of disruption threshold $\delta$ on concept removal and damage rates ($64\times$ SAE). At the lenient $\delta{>}2\%$ threshold, many features register as ``damaged'' from minor perturbation; at $\delta{>}20\%$, only features with substantial activation change count. TanGCE's damage drops to 1--11\% at $\delta{>}20\%$, while ConceptSAE Soft retains 17--80\% damage, confirming its lack of surgical precision.}
\label{tab:damage_threshold}
\small
\setlength{\tabcolsep}{2.5pt}
\begin{tabular}{l|cc|cc|cc}
\toprule
\multicolumn{7}{c}{\textbf{Sycophancy}} \\
 & \multicolumn{2}{c|}{$\delta{>}2\%$} & \multicolumn{2}{c|}{$\delta{>}10\%$} & \multicolumn{2}{c}{$\delta{>}20\%$} \\
Method & Rem\% & Dmg\% & Rem\% & Dmg\% & Rem\% & Dmg\% \\
\midrule
Random Global & 0 & 5 & 0 & 0 & 0 & 0 \\
Linear Projection & 0 & 24 & 0 & 5 & 0 & 3 \\
INLP & 0 & 24 & 0 & 5 & 0 & 3 \\
LEACE & 0 & 3 & 0 & 0 & 0 & 0 \\
Ambient GCE & 0 & 3 & 0 & 0 & 0 & 0 \\
TanGCE (raw) & 33 & 61 & 0 & 18 & 0 & 11 \\
ConceptSAE Soft & 50 & 70 & 0 & 25 & 0 & 17 \\
\midrule
\midrule
\multicolumn{7}{c}{\textbf{Gender}} \\
 & \multicolumn{2}{c|}{$\delta{>}2\%$} & \multicolumn{2}{c|}{$\delta{>}10\%$} & \multicolumn{2}{c}{$\delta{>}20\%$} \\
Method & Rem\% & Dmg\% & Rem\% & Dmg\% & Rem\% & Dmg\% \\
\midrule
Random Global & 50 & 26 & 0 & 1 & 0 & 0 \\
Linear Projection & 0 & 29 & 0 & 0 & 0 & 0 \\
INLP & 0 & 28 & 0 & 0 & 0 & 0 \\
LEACE & 100 & 35 & 0 & 1 & 0 & 0 \\
Ambient GCE & 0 & 22 & 0 & 0 & 0 & 0 \\
TanGCE (raw) & 50 & 61 & 0 & 10 & 0 & 1 \\
ConceptSAE Soft & 100 & 97 & 100 & 90 & 50 & 80 \\
\midrule
\midrule
\multicolumn{7}{c}{\textbf{Safety}} \\
 & \multicolumn{2}{c|}{$\delta{>}2\%$} & \multicolumn{2}{c|}{$\delta{>}10\%$} & \multicolumn{2}{c}{$\delta{>}20\%$} \\
Method & Rem\% & Dmg\% & Rem\% & Dmg\% & Rem\% & Dmg\% \\
\midrule
Random Global & 13 & 13 & 0 & 0 & 0 & 0 \\
Linear Projection & 0 & 2 & 0 & 0 & 0 & 0 \\
INLP & 0 & 2 & 0 & 0 & 0 & 0 \\
LEACE & 4 & 7 & 0 & 0 & 0 & 0 \\
Ambient GCE & 4 & 3 & 0 & 0 & 0 & 0 \\
TanGCE (raw) & 83 & 64 & 26 & 21 & 13 & 7 \\
ConceptSAE Soft & 87 & 75 & 43 & 34 & 30 & 25 \\
\bottomrule
\end{tabular}
\end{table}

Table~\ref{tab:damage_threshold} reveals that the high damage rates at $\delta{>}2\%$ are largely an artifact of the lenient threshold: many non-concept features shift by 2--5\% in mean activation---a perturbation too small to meaningfully alter downstream behavior (consistent with preserved MI$_\text{ctrl}$ in Table~\ref{tab:full_comparison_detail}). At the strict $\delta{>}20\%$ threshold, TanGCE's damage drops to 1\% on gender and 7\% on safety, confirming that TanGCE's apparent collateral impact at the lenient threshold reflects superficial activation shifts, not genuine information disruption.

\subsection{The $\lambda$ Sweep: Off-Manifold Displacement Creates Spurious Correlations}
\label{sec:lambda_sweep}

\begin{table}[t]
\centering
\caption{$\lambda$ sweep comparing \AmbGCE{} and \TanGCE{} at fixed step sizes across three concept domains (Qwen2.5-0.5B, layer~8). NL: retrained nonlinear MLP probe accuracy (hidden 256, 1500 steps, AdamW lr=$10^{-3}$; same architecture used for HPO and evaluation) on the target concept (lower = stronger erasure). Both methods use the same GCE concept scorer (hidden 256, 1500 steps). \AmbGCE{} becomes less effective at high $\lambda$ under the current protocol, while \TanGCE{} continues improving. We interpret this as supportive evidence for the local manifold motivation of tangent-constrained updates, not as a complete characterization of the leakage-utility frontier. The expanded HPO grid $\{0.25,\ldots,32.0\}$ used in the multi-model experiments further tests this range. Bold: best NL per domain.}
\label{tab:lambda_sweep}
\small
\setlength{\tabcolsep}{5pt}
\begin{tabular}{ll|ccc}
\toprule
 & & \multicolumn{3}{c}{NL$\downarrow$ at $\lambda=$} \\
Concept & Method & $0.5$ & $4.0$ & $8.0$ \\
\midrule
\multirow{2}{*}{Sycophancy}
 & \AmbGCE{} & .678 & .694 & .734 \\
 & \TanGCE{}  & .662 & .596 & \textbf{.588} \\
\midrule
\multirow{2}{*}{Gender}
 & \AmbGCE{} & .666 & .700 & .702 \\
 & \TanGCE{}  & .704 & .580 & \textbf{.470} \\
\midrule
\multirow{2}{*}{Safety}
 & \AmbGCE{} & .664 & .776 & .770 \\
 & \TanGCE{}  & .748 & .680 & \textbf{.646} \\
\bottomrule
\end{tabular}
\end{table}


Table~\ref{tab:lambda_sweep} reveals a critical failure mode of Ambient GCE: at $\lambda \geq 4$, concept MI \emph{increases} rather than decreases (negative MI drop). On sycophancy at $\lambda{=}8$, the MI drop is $-0.063$; on gender, $-0.098$; on safety, $-0.092$. Simultaneously, redistribution reaches 15--31\%.

This occurs because ambient erasure pushes representations off the data manifold. The displacement magnitude correlates with each sample's concept signal strength, creating a new axis of variation that SAE features detect as concept-correlated. Concept information is laundered into the geometry of the displacement rather than removed.

TanGCE avoids this failure mode through two mechanisms. First, projection onto the tangent space removes off-manifold gradient components, ensuring displacement stays on the manifold. Second, retaining the raw projection norm $\|\mathbf{P}_T(\nabla f)\|$ as a confidence weight naturally attenuates erasure for samples where the concept gradient is mostly off-manifold. This yields monotonically decreasing leakage with increasing $\lambda$, with redistribution remaining below 8\% even at $\lambda{=}8$.

\subsection{Polysemanticity Confounds Linear Correlation Metrics}
\label{sec:polysemanticity}

\begin{table}[t]
\centering
\caption{Effect of SAE dictionary size on concept removal detection. We train TopK SAEs ($k{=}64$) at $8\times$ and $64\times$ expansion on the same general-purpose WikiText corpus and re-evaluate all methods. Concept removal rate (\%) measures how many concept-correlated features are disrupted by erasure; damage rate (\%) measures collateral disruption to non-concept features. Correlation threshold $|r|{>}0.1$; disruption threshold $\delta{>}0.02$. The $8\times$ SAE (7{,}168 features) systematically underestimates concept removal due to polysemantic feature representations, particularly on gender where it reports 0\% removal for methods that demonstrably erase the concept.}
\label{tab:sae_expansion}
\small
\setlength{\tabcolsep}{3pt}
\begin{tabular}{l|cc|cc|cc|cc|cc|cc}
\toprule
 & \multicolumn{4}{c|}{Sycophancy} & \multicolumn{4}{c|}{Gender (Bias in Bios)} & \multicolumn{4}{c}{Safety} \\
 & \multicolumn{2}{c|}{Remove\%} & \multicolumn{2}{c|}{Damage\%} & \multicolumn{2}{c|}{Remove\%} & \multicolumn{2}{c|}{Damage\%} & \multicolumn{2}{c|}{Remove\%} & \multicolumn{2}{c}{Damage\%} \\
Method & $8\times$ & $64\times$ & $8\times$ & $64\times$ & $8\times$ & $64\times$ & $8\times$ & $64\times$ & $8\times$ & $64\times$ & $8\times$ & $64\times$ \\
\midrule
Linear Proj. & 14 & 14 & 14 & 19 & 0 & 0 & 23 & 29 & 0 & 0 & 0 & 2 \\
INLP         & 14 & 14 & 14 & 19 & 0 & 0 & 23 & 28 & 0 & 0 & 0 & 2 \\
LEACE        & 14 & 0  & 1  & 3  & \textbf{0} & \textbf{100} & 29 & 35 & 9 & 5 & 4 & 5 \\
Ambient GCE  & 0  & 0  & 3  & 1  & 0 & 0 & 16 & 20 & 0 & 5 & 0 & 2 \\
TanGCE (raw) & 71 & 71 & 61 & 60 & \textbf{0} & \textbf{100} & 51 & 60 & 59 & 82 & 65 & 68 \\
ConceptSAE   & 57 & 71 & 76 & 63 & 0 & 100 & 99 & 98 & 69 & 91 & 76 & 76 \\
\bottomrule
\end{tabular}
\end{table}


A natural approach for measuring concept removal in SAE feature space is to identify concept-correlated features via point-biserial correlation $r(z_j, y)$ and then measure what fraction are disrupted after erasure. However, this approach assumes that each SAE feature encodes a single semantic concept---an assumption that fails when features are \emph{polysemantic}.

Consider a polysemantic feature $z_j$ that activates when \emph{both} concept $y{=}1$ and topic $t{=}\text{science}$ are present, or when $y{=}0$ and $t{=}\text{politics}$. The marginal correlation $r(z_j, y) \approx 0$ because the feature fires equally often for both concept labels, yet $z_j$ carries substantial concept information \emph{conditioned on topic}. A linear correlation metric would classify this feature as non-concept-related and miss the fact that erasure correctly disrupts it.

Table~\ref{tab:sae_expansion} demonstrates this effect empirically. At $8\times$ expansion (7{,}168 dictionary features), the SAE reports \textbf{0\% concept removal} for both TanGCE and LEACE on gender---despite both methods substantially reducing nonlinear probe leakage (Table~\ref{tab:full_comparison}). Increasing the dictionary to $64\times$ expansion (57{,}344 features) resolves the polysemanticity: features become more monosemantic, linear correlations faithfully detect concept encoding, and both methods now show \textbf{100\% concept removal}.

The effect is concept-dependent. Sycophancy removal rates are stable across expansions (71\% at both $8\times$ and $64\times$ for TanGCE), suggesting that sycophancy features are already relatively monosemantic at $8\times$. Gender features, by contrast, are heavily polysemantic at $8\times$---the concept is distributed across features that jointly encode gender with other attributes, making it invisible to marginal correlation. Safety shows an intermediate pattern: removal increases from 59\% to 82\%.

These results have two implications. First, SAE-based erasure evaluation requires sufficiently high dictionary expansion to overcome polysemanticity; $8\times$ expansion is inadequate for concepts that are entangled with other attributes. Second, the zero-redistribution result for TanGCE on gender (Table~\ref{tab:full_comparison}) is robust: redistribution remains at 0\% even at $64\times$ expansion, confirming that TanGCE removes rather than reshuffles concept information.

\subsection{Geometric Evidence and Off-Manifold Displacement}
\begin{table}[t]
\centering
\caption{Geometry contrast between sycophancy and gender across completed layers 4, 8, and 12. This is a focused contrast between the primary nonlinear case and the diagnostic baseline, not an all-dataset geometry census. The table reports two quantities: a local intrinsic-dimension range and the mean angle between a global concept direction and local per-sample concept directions. Under the manifold framing, large local-global angles indicate that a single global direction is a poor first-order summary across neighborhoods. They do not by themselves provide a direct estimate of curvature. Sycophancy shows substantially larger local-global misalignment than gender, supporting it as the primary case for locally adaptive, manifold-motivated erasure methods.}
\label{tab:geometry_contrast_summary}
\small
\begin{tabular}{lcc}
\toprule
Concept & Local intrinsic dimension & Mean local/global angle \\
\midrule
Sycophancy & 34--53 & 76.9--82.6$^\circ$ \\
Gender & 365--461 & 51.8--55.8$^\circ$ \\
\bottomrule
\end{tabular}
\end{table}


Table~\ref{tab:geometry_contrast_summary} demonstrates that sycophancy is much more nonlinear than gender: its intrinsic dimension (34--53 across layers 4, 8, 12) is far lower than gender's (365--461), and its local-vs-global angle is much larger (76.9--82.6$^\circ$ vs.\ 51.8--55.8$^\circ$). This geometry contrast validates sycophancy as the primary case for manifold-aware erasure methods, while the strong TanGCE results on gender demonstrate that the method generalizes even to approximately linear concepts.

\paragraph{Appendix-Supporting Evidence.}
Detailed geometric analyses supporting these claims are provided in the Appendix, including:
\begin{itemize}
    \item \textbf{Subspace Sufficiency (Q1):} Evidence that task-specific computation is concentrated in a low-dimensional linear subspace.
    \item \textbf{Perturbation Sensitivity (Q2):} Comparison of perturbations in high-variance and orthogonal subspaces.
    \item \textbf{Representation Connectivity (Q3):} Analysis of the connected, approximately convex topology of the representation manifold.
    \item \textbf{Multi-Layer Accumulation:} Results on how projection error accumulates across the network.
\end{itemize}

% Discussion
\section{Discussion}
Content will be populated from PAPER.md/reports in later tasks.

% Limitations
\section{Limitations}

\paragraph{Single-seed snapshot.} The headline tables report a single deterministic seed=0 run on the 39 NLP settings and the 80 CelebA-CLIP settings. Multi-seed variance across the 39 NLP settings will be reported in App.\ as the additional runs complete; the seed=0 numbers are not preregistered as the final headline values, only the result available at submission.

\paragraph{Measurement scope.} Concept removal is judged by retraining a fresh 2-layer MLP probe ($h{=}128$, 200 SGD steps, patience 3) on the edited representations. This captures recoverability under the protocol of \citet{elazar-goldberg-2018-adversarial} and surgicality through control-task accuracy on the labelled control concepts. It does not by itself establish broad capability preservation across unrelated tasks: a corpus-level evaluation on held-out general-domain text is left to future work. The cross-modal evidence comes from a single CLIP encoder (ViT-B/32 on CelebA); evaluations against other vision encoders or domains are out of scope here.

\paragraph{Geometry as supporting evidence, not causal proof.} The local manifold motivation is operational, not a theorem. TwoNN intrinsic-dimension estimates and secant-SVD bases are first-order approximations of an unobserved tangent space. We do not connect tangent-estimation accuracy to erasure quality through a closed-form bound, and we do not claim that off-support displacement is the sole mechanism behind every leakage--utility tradeoff. The geometry diagnostics support tangent projection as a useful inductive bias rather than as a complete causal account.

% Conclusion
\section{Conclusion}
Content will be populated from PAPER.md/reports in later tasks.


\bibliographystyle{plainnat}
\bibliography{references}

% Keep appendix marker (actual appendix file will be created in follow-up task)
\appendix

% Appendix snapshot tables
\section{Repository Snapshot Tables}
\label{sec:appendix_snapshots}

This appendix records broader repository snapshot tables that are directionally consistent with the main text but were not part of the locked paper-evidence set when the quantitative claims in Section~\ref{sec:results} were written. We include them to document the current state of the project and to guide the next round of matched-leakage and deeper-layer analyses. The central paper claims continue to rely only on the completed paper suite summarized in the main text.

\subsection{Iterative Sycophancy Snapshot}
\begin{table}[t]
\centering
\caption{Repository snapshot from the iterative sycophancy report at $\lambda=1.0$ after 7 rounds, across layers 8, 16, and 23. Columns T (tangent-constrained) and A (ambient baseline) compare erasure methods. Utility is answer-matching delta relative to clean representations. Leakage and displacement decrease from left to right (lower is better); utility increases from left to right (higher is better). This snapshot table is not a main-text claim anchor and should be read only as evidence of directional consistency with the layer-8 paper-locked row (Table~\ref{tab:completed_results_summary}). Deeper-layer iterative tangent runs continue to reduce leakage more than ambient baselines across this exploratory sweep.}
\label{tab:repo_snapshot_sycophancy}
\small
\begin{tabular}{lcccccc}
\toprule
Layer & T leak & A leak & T util & A util & T disp & A disp \\
\midrule
8 & 0.063 & 0.213 & -20.7pp & -0.5pp & 2.88 & 0.62 \\
16 & 0.153 & 0.400 & +0.0pp & +0.0pp & 3.79 & 0.27 \\
23 & 0.138 & 0.361 & -0.1pp & -0.2pp & 44.50 & 3.11 \\
\bottomrule
\end{tabular}
\end{table}


\subsection{Iterative Safety Snapshot}
\begin{table}[t]
\centering
\caption{Repository snapshot from the iterative safety report at $\lambda=1.0$ after 7 rounds. Utility is pairwise helpfulness delta. These rows are directionally consistent with the main text but are not used as main-text claim anchors.}
\label{tab:repo_snapshot_safety}
\small
\begin{tabular}{lcccccc}
\toprule
Layer & T leak & A leak & T util & A util & T disp & A disp \\
\midrule
8 & 0.0725 & 0.2583 & -0.052 & -0.023 & 2.70 & 0.18 \\
16 & 0.1292 & 0.3517 & -0.092 & +0.003 & 4.37 & 0.60 \\
23 & 0.1483 & 0.3404 & -0.062 & -0.152 & 38.32 & 5.30 \\
\bottomrule
\end{tabular}
\end{table}


\subsection{Local CelebA Visual Table-2-Style Comparison}
\begin{table*}[t]
\centering
\caption{Local supplementary Table-2-style CelebA visual comparison using CLIP ViT-B/32 embeddings. For each target concept we intervene once per method and evaluate (i) concept leakage via retrained linear/nonlinear probes (lower is better) and (ii) mean control-task accuracy change across 3 least-correlated control attributes (closer to zero is better). Targets and controls: Male $\rightarrow$ Blurry,Narrow\_Eyes,Straight\_Hair; Smiling $\rightarrow$ Straight\_Hair,Bushy\_Eyebrows,Gray\_Hair; Blond\_Hair $\rightarrow$ Straight\_Hair,Narrow\_Eyes,Blurry. In this local supplementary setting, all methods use $\lambda{=}0.5$ except \TanGCE{} with $\lambda{=}8.0$.}
\label{tab:local_celeba_visual_table2}
\small
\setlength{\tabcolsep}{4pt}
\begin{tabular}{l|ccc|ccc|ccc}
\toprule
 & \multicolumn{3}{c|}{Male} & \multicolumn{3}{c|}{Smiling} & \multicolumn{3}{c}{Blond\_Hair} \\
Method & L$\downarrow$ & NL$\downarrow$ & $\Delta$Ctrl & L$\downarrow$ & NL$\downarrow$ & $\Delta$Ctrl & L$\downarrow$ & NL$\downarrow$ & $\Delta$Ctrl \\
\midrule
Random Global & .993 & .994 & -0.0pp & .891 & .911 & +0.0pp & .909 & .933 & -0.1pp \\
Random Per-Samp. & .993 & .994 & +0.0pp & .891 & .907 & +0.2pp & .907 & .932 & -0.1pp \\
Linear Proj. & .993 & .992 & +0.0pp & .900 & .906 & -0.3pp & .919 & .931 & -0.2pp \\
INLP & .994 & .993 & -0.0pp & .900 & .903 & +0.3pp & .918 & .932 & -0.0pp \\
LEACE & .994 & .991 & -0.0pp & .894 & .894 & +0.2pp & .909 & .925 & -0.1pp \\
\AmbGCE{} & .990 & .995 & -0.2pp & .858 & .884 & -0.4pp & .893 & .920 & +0.2pp \\
\TanGCE{} & .887 & .953 & -0.8pp & .649 & .649 & -0.6pp & .699 & .747 & -2.1pp \\
\bottomrule
\end{tabular}
\end{table*}


To ensure a modality-matched benchmark, this appendix now reports a vision-only CelebA run (CLIP ViT-B/32 embeddings) with the Table~\ref{tab:full_comparison} method family and 3 least-correlated controls per target concept. Artifacts are stored under \texttt{outputs/local/celeba\_visual\_table2\_v1/} (including per-method sample folders with 2 saved images each). The consolidated W\&B results run is \texttt{g52h8kbr}.

\subsection{CelebA Per-Concept Detail (All 40 Attributes)}
\label{sec:appendix_celeba_per_concept}
Tables~\ref{tab:celeba_clip_all40_per_concept_leak} and~\ref{tab:celeba_clip_all40_per_concept_ctrl} give the full per-concept breakdown behind the aggregate CelebA result in Table~\ref{tab:celeba_clip_all40}. The main-text Table~\ref{tab:celeba_clip_fair} reports three representative concepts (Male, Smiling, Blond\_Hair); the appendix tables cover all 40 binary CelebA attributes from the same \texttt{celeba\_clip\_all40\_fair\_v2} run under the shared validation-HPO protocol. TanGCE achieves the lowest nonlinear leakage on \textbf{36 of 40} concepts (LEACE wins on 3, IGBP on 1), with a mean margin of $0.116$ absolute NL accuracy over the best competing method per concept. Per-concept mean $\Delta$Ctrl remains within roughly a percentage point of zero, consistent with the aggregate mean NL of $0.707$ and mean $\Delta$Ctrl of $-0.7\%$ reported in the main text.
\begin{table*}[p]
\centering
\caption{Per-concept nonlinear probe leakage on CelebA (CLIP ViT-B/32), all 40 binary attributes. Each cell is retrained nonlinear probe test accuracy after erasure (lower = stronger erasure). $\lambda^*$ is selected per method per concept from $\{0.5, 1.0, 2.0, 4.0, 8.0\}$ by lowest validation NL accuracy. Bold marks the lowest leakage per row. Abbrev.: RandG = Random Global, RandS = Random Per-Sample, LinProj = Linear Projection, \AmbGCE{} = ambient gradient-based concept erasure, \TanGCE{} = tangent gradient-based concept erasure.}
\label{tab:celeba_clip_all40_per_concept_leak}
\footnotesize
\setlength{\tabcolsep}{3pt}
\begin{tabular}{l|cccccccc}
\toprule
Concept & RandG & RandS & LinProj & INLP & LEACE & IGBP & \AmbGCE{} & \TanGCE{} \\
\midrule
5\_o\_Clock\_Shadow & .92 & .92 & .91 & .92 & .87 & .92 & .91 & \textbf{.77} \\
Arched\_Eyebrows & .79 & .79 & .76 & .76 & .69 & .77 & .77 & \textbf{.59} \\
Attractive & .73 & .74 & .69 & .73 & .63 & \textbf{.53} & .73 & .60 \\
Bags\_Under\_Eyes & .76 & .75 & .74 & .75 & .70 & .77 & .75 & \textbf{.69} \\
Bald & 1.00 & .99 & .99 & .99 & .98 & 1.00 & .99 & \textbf{.97} \\
Bangs & .91 & .91 & .89 & .89 & .87 & .92 & .91 & \textbf{.68} \\
Big\_Lips & .80 & .79 & .74 & .80 & \textbf{.68} & .82 & .81 & .73 \\
Big\_Nose & .76 & .76 & .73 & .74 & .72 & .73 & .76 & \textbf{.61} \\
Black\_Hair & .85 & .84 & .82 & .81 & .79 & .83 & .82 & \textbf{.46} \\
Blond\_Hair & .92 & .92 & .94 & .92 & .92 & .92 & .92 & \textbf{.72} \\
Blurry & .98 & .98 & .98 & .96 & .97 & .98 & .97 & \textbf{.88} \\
Brown\_Hair & .80 & .80 & .76 & .75 & .75 & .80 & .80 & \textbf{.52} \\
Bushy\_Eyebrows & .87 & .88 & .86 & .85 & .76 & .87 & .86 & \textbf{.66} \\
Chubby & .94 & .94 & .95 & .95 & .92 & .95 & .94 & \textbf{.81} \\
Double\_Chin & .94 & .94 & .94 & .93 & .92 & .94 & .95 & \textbf{.85} \\
Eyeglasses & .99 & 1.00 & .96 & .93 & .96 & 1.00 & .96 & \textbf{.90} \\
Goatee & .96 & .97 & .97 & .93 & .96 & .96 & .95 & \textbf{.85} \\
Gray\_Hair & .97 & .97 & .97 & .97 & .97 & .97 & .98 & \textbf{.89} \\
Heavy\_Makeup & .89 & .89 & .89 & .86 & .85 & .88 & .88 & \textbf{.70} \\
High\_Cheekbones & .79 & .79 & .77 & .75 & .70 & .57 & .76 & \textbf{.50} \\
Male & .99 & .99 & .99 & .99 & 1.00 & .98 & 1.00 & \textbf{.85} \\
Mouth\_Slightly\_Open & .89 & .89 & .89 & .71 & .72 & .66 & .80 & \textbf{.43} \\
Mustache & .94 & .95 & .95 & .94 & .92 & .95 & .93 & \textbf{.77} \\
Narrow\_Eyes & .88 & .88 & .80 & .87 & .80 & .88 & .88 & \textbf{.78} \\
No\_Beard & .94 & .95 & .94 & .93 & .92 & .94 & .95 & \textbf{.77} \\
Oval\_Face & .70 & .71 & .69 & .68 & \textbf{.67} & .69 & .67 & .72 \\
Pale\_Skin & .93 & .93 & .94 & .95 & .93 & .92 & .93 & \textbf{.83} \\
Pointy\_Nose & .75 & .72 & .71 & .69 & \textbf{.63} & .72 & .75 & .64 \\
Receding\_Hairline & .94 & .94 & .96 & .96 & .94 & .94 & .95 & \textbf{.69} \\
Rosy\_Cheeks & .93 & .93 & .93 & .94 & .91 & .94 & .93 & \textbf{.79} \\
Sideburns & .95 & .95 & .94 & .94 & .95 & .94 & .94 & \textbf{.82} \\
Smiling & .90 & .89 & .91 & .74 & .82 & .59 & .88 & \textbf{.44} \\
Straight\_Hair & .78 & .78 & .79 & .77 & .78 & .78 & .78 & \textbf{.49} \\
Wavy\_Hair & .78 & .76 & .77 & .77 & .78 & .77 & .78 & \textbf{.52} \\
Wearing\_Earrings & .90 & .89 & .86 & .88 & .85 & .89 & .88 & \textbf{.61} \\
Wearing\_Hat & 1.00 & 1.00 & .97 & .97 & .97 & 1.00 & .97 & \textbf{.94} \\
Wearing\_Lipstick & .90 & .90 & .90 & .89 & .85 & .88 & .92 & \textbf{.75} \\
Wearing\_Necklace & .84 & .85 & .80 & .80 & .77 & .85 & .83 & \textbf{.68} \\
Wearing\_Necktie & .93 & .93 & .92 & .92 & .91 & .92 & .93 & \textbf{.78} \\
Young & .86 & .86 & .85 & .84 & .83 & .86 & .87 & \textbf{.60} \\
\midrule
\textbf{Mean} & .88 & .88 & .87 & .86 & .84 & .86 & .87 & \textbf{.71} \\
\bottomrule
\end{tabular}
\end{table*}

\begin{table*}[p]
\centering
\caption{Per-concept mean control accuracy change $\Delta$Ctrl (\%) on CelebA, all 40 binary attributes. For each concept, $\Delta$Ctrl is averaged over the 5 least-correlated control attributes (closer to zero = more surgical). $\lambda^*$ selection matches Table~\ref{tab:celeba_clip_all40_per_concept_leak}.}
\label{tab:celeba_clip_all40_per_concept_ctrl}
\footnotesize
\setlength{\tabcolsep}{3pt}
\begin{tabular}{l|cccccccc}
\toprule
Concept & RandG & RandS & LinProj & INLP & LEACE & IGBP & \AmbGCE{} & \TanGCE{} \\
\midrule
5\_o\_Clock\_Shadow & +0.0 & +0.1 & -0.7 & +0.1 & +0.0 & +0.1 & -2.3 & -0.2 \\
Arched\_Eyebrows & -1.3 & -2.3 & -1.2 & -0.8 & -1.1 & -2.5 & +0.6 & -0.6 \\
Attractive & -0.3 & -1.0 & -1.7 & -0.8 & -0.3 & -0.9 & +0.0 & -0.8 \\
Bags\_Under\_Eyes & -3.7 & -1.7 & -0.3 & -7.7 & -0.3 & -0.3 & -0.3 & -0.1 \\
Bald & +0.1 & +0.0 & -0.1 & +0.1 & -0.1 & +1.1 & +2.3 & +0.2 \\
Bangs & -3.7 & -1.2 & -0.5 & -7.9 & -0.1 & -0.7 & -0.3 & -0.8 \\
Big\_Lips & +0.1 & -0.0 & -0.3 & +0.0 & -0.1 & -0.5 & +0.1 & -0.1 \\
Big\_Nose & -0.1 & +0.3 & -3.5 & -9.7 & +0.0 & -0.3 & -1.8 & +0.1 \\
Black\_Hair & -0.1 & -0.0 & +0.1 & +0.2 & +0.3 & -0.3 & -0.3 & -0.7 \\
Blond\_Hair & -0.3 & +0.1 & -1.4 & +0.7 & -0.1 & -0.4 & -1.1 & -1.6 \\
Blurry & -1.3 & +0.0 & -0.1 & -5.5 & +0.1 & -0.7 & +0.7 & -0.1 \\
Brown\_Hair & +0.1 & -0.9 & -0.3 & -6.0 & -0.7 & +0.4 & +0.1 & -0.2 \\
Bushy\_Eyebrows & +0.0 & +0.1 & -0.1 & -1.5 & +0.2 & -0.9 & -0.3 & -0.5 \\
Chubby & -0.1 & +0.3 & -0.2 & -1.4 & -0.2 & -0.2 & +0.4 & -2.0 \\
Double\_Chin & +0.1 & +0.1 & -0.4 & -12.7 & +0.3 & +0.7 & -0.7 & -2.5 \\
Eyeglasses & -0.1 & +0.1 & -0.0 & -0.1 & -0.2 & -0.4 & -0.2 & +0.2 \\
Goatee & +0.1 & -0.1 & +0.0 & +0.2 & +0.1 & +0.5 & -0.3 & -0.4 \\
Gray\_Hair & -0.2 & -0.3 & -2.1 & +0.4 & -0.1 & -1.1 & -1.5 & -1.5 \\
Heavy\_Makeup & -0.3 & -0.1 & -1.1 & -1.7 & -1.3 & -1.3 & -1.1 & -0.9 \\
High\_Cheekbones & +0.1 & +0.2 & +0.1 & -0.3 & -0.1 & -0.4 & -0.0 & -0.7 \\
Male & -0.3 & -0.1 & -0.7 & -0.2 & -0.3 & -2.2 & -0.9 & -1.1 \\
Mouth\_Slightly\_Open & -0.3 & -0.6 & +0.1 & +0.4 & +0.2 & +0.3 & +0.1 & -0.2 \\
Mustache & +0.0 & -0.1 & +0.1 & -10.3 & -0.5 & +1.0 & +0.5 & -1.5 \\
Narrow\_Eyes & -0.1 & -0.2 & -0.3 & +0.2 & -0.1 & +0.3 & +0.0 & +0.0 \\
No\_Beard & -1.2 & -0.2 & -4.0 & -6.5 & -1.6 & -0.1 & -0.7 & -0.5 \\
Oval\_Face & -2.0 & +0.1 & -0.1 & -14.4 & -0.6 & -0.3 & +0.2 & -0.3 \\
Pale\_Skin & -1.2 & -0.1 & -0.1 & -6.8 & -0.9 & -0.1 & +0.7 & -1.0 \\
Pointy\_Nose & -0.1 & -0.2 & +0.3 & -1.1 & +0.2 & -0.7 & -0.1 & -0.1 \\
Receding\_Hairline & -0.1 & -0.1 & -0.1 & -5.9 & +0.0 & -0.3 & +0.1 & -1.1 \\
Rosy\_Cheeks & -0.5 & +0.1 & -0.6 & -0.1 & -0.2 & -0.1 & -0.1 & -1.3 \\
Sideburns & -0.4 & +0.0 & -0.1 & +0.1 & +0.1 & -0.6 & -2.5 & -0.1 \\
Smiling & -0.3 & -1.3 & +0.1 & +0.6 & +0.4 & +0.1 & +0.1 & -0.1 \\
Straight\_Hair & -0.1 & -0.2 & -0.1 & -0.9 & -0.1 & -0.8 & +0.5 & -0.7 \\
Wavy\_Hair & +0.1 & -0.7 & -0.9 & +0.2 & +0.1 & -0.4 & -0.1 & -0.7 \\
Wearing\_Earrings & +0.8 & -0.3 & -1.6 & -0.3 & -0.1 & -1.3 & -0.5 & -0.3 \\
Wearing\_Hat & -0.1 & +0.0 & -0.1 & -0.6 & +0.0 & +1.0 & +0.1 & -0.2 \\
Wearing\_Lipstick & +0.2 & -0.1 & -6.1 & -21.5 & -1.9 & -1.2 & -1.4 & -0.9 \\
Wearing\_Necklace & +0.1 & -1.7 & -1.7 & -9.3 & +0.3 & -1.8 & +0.2 & -1.2 \\
Wearing\_Necktie & +0.0 & +0.1 & +0.0 & +0.2 & -0.2 & -0.7 & -13.4 & -1.3 \\
Young & -0.0 & -0.7 & -1.7 & -10.7 & +0.1 & -0.9 & -2.2 & -0.8 \\
\midrule
\textbf{Mean} & -0.4 & -0.3 & -0.8 & -3.5 & -0.2 & -0.4 & -0.6 & -0.7 \\
\bottomrule
\end{tabular}
\end{table*}


\begin{table}[t]
\centering
\caption{Geometry evidence map organizing the appendix diagnostics into four research blocks. Each block addresses a distinct aspect of manifold structure: subspace sufficiency examines whether task-relevant computation concentrates in a low-dimensional linear subspace, perturbation sensitivity tests whether on-manifold directions matter more than orthogonal directions, representation connectivity evaluates whether the manifold remains connected across domains, and multi-layer accumulation tracks error compounding through the network. These diagnostics support the mechanism discussion and motivate the need for manifold-aware erasure methods but do not by themselves establish causality.}
\label{tab:geometry_schedule}
\scriptsize
\begin{tabular}{p{0.22\linewidth}p{0.22\linewidth}p{0.40\linewidth}}
\toprule
Experiment Block & Research Question & Primary Evidence \\
\midrule
Subspace sufficiency & Subspace concentration & Top-1 agreement under intrinsic projection \\
Perturbation sensitivity & High-variance vs. orthogonal directions & KL and top-1 under on/off-subspace noise \\
Representation connectivity & Shared vs. domain manifolds & Cross-domain interpolation smoothness \\
Multi-layer accumulation & Error compounding across the network & Simultaneous intrinsic projection across layers \\
\bottomrule
\end{tabular}
\end{table}


\paragraph{Geometry Note.}
The broader geometry report continues to support three useful points: early and mid-layer representations are effectively low-dimensional, on-manifold directions carry much of the model's sensitivity in these perturbation tests, and interpolation remains consistent with a connected representation space. At the same time, the Q2 perturbation study is best read as evidence of sensitivity concentration in a high-variance subspace rather than as a standalone proof of curved-manifold causality.


\newpage
\input{checklist.tex}

\end{document}
